\section{AMC 12 12A 2008}
        \begin{enumerate}
        	\item (\textit{AMC 12 12A 2008}) A bakery owner turns on his doughnut machine at $\text{8:30}\ {\small\text{AM}}$. At $\text{11:10}\ {\small\text{AM}}$ the machine has completed one third of the day's job. At what time will the doughnut machine complete the job?


 $\mathrm{(A)}\ \text{1:50}\ {\small\text{PM}}\qquad\mathrm{(B)}\ \text{3:00}\ {\small\text{PM}}\qquad\mathrm{(C)}\ \text{3:30}\ {\small\text{PM}}\qquad\mathrm{(D)}\ \text{4:30}\ {\small\text{PM}}\qquad\mathrm{(E)}\ \text{5:50}\ {\small\text{PM}}$

 \textbf{Answer: }D



\textbf{Solution}

The machine completes one-third of the job in $\text{11:10}-\text{8:30}=\text{2:40}$ hours. Thus, the entire job is completed in $3\cdot(\text{2:40})=\text{8:00}$ hours.

 Since the machine was started at $\text{8:30 AM}$, the job will be finished $8$ hours later, at $\text{4:30 PM}$. The answer is $\mathrm{(D)}$.

 Note: $\text{2:40}$ means $2$ hours and $40$ minutes. $3$ multiplied by this time interval is $8$ hours.

 


	\item (\textit{AMC 12 12A 2008}) What is the \href{https://artofproblemsolving.com/wiki/index.php?title=Reciprocal}{reciprocal} of $\frac{1}{2}+\frac{2}{3}$?


 $\mathrm{(A)}\ \frac{6}{7}\qquad\mathrm{(B)}\ \frac{7}{6}\qquad\mathrm{(C)}\ \frac{5}{3}\qquad\mathrm{(D)}\ 3\qquad\mathrm{(E)}\ \frac{7}{2}$

 \textbf{Answer: }A



\textbf{Solution 1}

Here's a cheapshot: Obviously, $\frac{1}{2}+\frac{2}{3}$ is greater than $1$. Therefore, its reciprocal is less than $1$, and the answer must be $\boxed{\frac{6}{7}}$.

 



\textbf{Solution 2}

$\left(\frac{1}{2}+\frac{2}{3}\right)^{-1}=\left(\frac{3}{6}+\frac{4}{6}\right)^{-1}=\left(\frac{7}{6}\right)^{-1}=\boxed{\mathrm{(A)}\ \frac{6}{7}}$.

 


	\item (\textit{AMC 12 12A 2008}) Suppose that $\frac{2}{3}$ of $10$ bananas are worth as much as $8$ oranges. How many oranges are worth as much as $\frac{1}{2}$ of $5$ bananas?


 $\mathrm{(A)}\ 2\qquad\mathrm{(B)}\ \frac{5}{2}\qquad\mathrm{(C)}\ 3\qquad\mathrm{(D)}\ \frac{7}{2}\qquad\mathrm{(E)}\ 4$

 \textbf{Answer: }C



\textbf{Solution}

If $\frac{2}{3}\cdot10\ \text{bananas}=8\ \text{oranges}$, then $\frac{1}{2}\cdot5\ \text{bananas}=\left(\frac{1}{2}\cdot 5\ \text{bananas}\right)\cdot\left(\frac{8\ \text{oranges}}{\frac{2}{3}\cdot10\ \text{bananas}}\right)=3\ \text{oranges}\Longrightarrow\mathrm{(C)}$. 

 


	\item (\textit{AMC 12 12A 2008}) Which of the following is equal to the \href{https://artofproblemsolving.com/wiki/index.php?title=Product}{product}

 \[\frac{8}{4}\cdot\frac{12}{8}\cdot\frac{16}{12}\cdot\cdots\cdot\frac{4n+4}{4n}\cdot\cdots\cdot\frac{2008}{2004}?\]
$\mathrm{(A)}\ 251\qquad\mathrm{(B)}\ 502\qquad\mathrm{(C)}\ 1004\qquad\mathrm{(D)}\ 2008\qquad\mathrm{(E)}\ 4016$

 \textbf{Answer: }B



\textbf{Solution 1}

$\frac {8}{4}\cdot\frac {12}{8}\cdot\frac {16}{12}\cdots\frac {4n + 4}{4n}\cdots\frac {2008}{2004} = \frac {1}{4}\cdot\left(\frac {8}{8}\cdot\frac {12}{12}\cdots\frac {4n}{4n}\cdots\frac {2004}{2004}\right)\cdot 2008 = \frac{2008}{4} =$ $502 \Rightarrow B$. 

 



\textbf{Solution 2}

Notice that everything cancels out except for $2008$ in the numerator and $4$ in the denominator.

 Thus, the product is $\frac{2008}{4}=502$, and the answer is $\textbf{(B)}$. 

 


	\item (\textit{AMC 12 12A 2008}) Suppose that

 \[\frac{2x}{3}-\frac{x}{6}\]
is an integer. Which of the following statements must be true about $x$?


 $\mathrm{(A)}\ \text{It is negative.}\qquad\mathrm{(B)}\ \text{It is even, but not necessarily a multiple of 3.}\\\qquad\mathrm{(C)}\ \text{It is a multiple of 3, but not necessarily even.}\\\qquad\mathrm{(D)}\ \text{It is a multiple of 6, but not necessarily a multiple of 12.}\\\qquad\mathrm{(E)}\ \text{It is a multiple of 12.}$

 \textbf{Answer: }B



\textbf{Solution}


 \[\frac{2x}{3}-\frac{x}{6}\quad\Longrightarrow\quad\frac{4x}{6}-\frac{x}{6}\quad\Longrightarrow\quad\frac{3x}{6}\quad\Longrightarrow\quad\frac{x}{2}\]For $\frac{x}{2}$ to be an integer, $x$ must be even, but not necessarily divisible by $3$. Thus, the answer is $\mathrm{(B)}$.

 


	\item (\textit{AMC 12 12A 2008}) Heather compares the price of a new computer at two different stores. Store $A$ offers $15\%$ off the sticker price followed by a $\textdollar 90$ rebate, and store $B$ offers $25\%$ off the same sticker price with no rebate. Heather saves $\textdollar 15$ by buying the computer at store $A$ instead of store $B$. What is the sticker price of the computer, in dollars?


 $\mathrm{(A)}\ 750\qquad\mathrm{(B)}\ 900\qquad\mathrm{(C)}\ 1000\qquad\mathrm{(D)}\ 1050\qquad\mathrm{(E)}\ 1500$

 \textbf{Answer: }A



\textbf{Solution 1}

Let the sticker price be $x$.

 The price of the computer is $0.85x-90$ at store $A$ and $0.75x$ at store $B$.

 Heather saves $\$15$ at store $A$, so $0.85x-90+15=0.75x$.

 Solving, we find $x=750$, and thus the answer is $\mathrm{(A)}$.

 



\textbf{Solution 2}

The $\$ 90$ at store $A$ is $\$ 15$ greater than the additional $25\%-15\% = 10\%$ off at store $B$.

 Thus the $10\%$ off is equal to $\$ 90$ $-$ $\$ 15$ $=$ $\$ 75$, and therefore the sticker price is $\$ 750$.

 


	\item (\textit{AMC 12 12A 2008}) While Steve and LeRoy are fishing 1 mile from shore, their boat springs a leak, and water comes in at a constant rate of 10 gallons per minute. The boat will sink if it takes in more than 30 gallons of water. Steve starts rowing towards the shore at a constant rate of 4 miles per hour while LeRoy bails water out of the boat. What is the slowest rate, in gallons per minute, at which LeRoy can bail if they are to reach the shore without sinking?


 $\mathrm{(A)}\ 2\qquad\mathrm{(B)}\ 4\qquad\mathrm{(C)}\ 6\qquad\mathrm{(D)}\ 8\qquad\mathrm{(E)}\ 10$

 \textbf{Answer: }D



\textbf{Solution 1}

It will take $\frac{1}{4}$ of an hour or $15$ minutes to get to shore.

 Since only $30$ gallons of water can enter the boat, only $\frac{30}{15}=2$ net gallons can enter the boat per minute.

 Since $10$ gallons of water enter the boat each minute, LeRoy must bail $10-2=8$ gallons per minute $\Rightarrow\mathrm{(D)}$.

 


 



\textbf{Solution 2}

We set up the following equation, where $x$ is the answer:

 $\frac{30}{10-x} = 15 \Rightarrow\mathrm{x=8}$, so the answer is $(D)$.

 



\textbf{Solution 3}

Let $r$ be the slowest rate LeRoy can bail out water. It will take them $15$ minutes to get to shore. As stated in solution 1, the number of gallons that enter the boat can be at most $2$ for the boat to stay afloat. Water enters at a rate of $10$ gallons per minute. 

 Therefore:

 $2 \geq 10-r$
$r \geq 8$

 The minimum possible value for $r$ is $8\ \mathrm{(D)}$.

 


	\item (\textit{AMC 12 12A 2008}) What is the \href{https://artofproblemsolving.com/wiki/index.php?title=Volume}{volume} of a \href{https://artofproblemsolving.com/wiki/index.php?title=Cube}{cube} whose \href{https://artofproblemsolving.com/wiki/index.php?title=Surface\_area}{surface area} is twice that of a cube with volume 1? 


 $\mathrm{(A)}\ \sqrt{2}\qquad\mathrm{(B)}\ 2\qquad\mathrm{(C)}\ 2\sqrt{2}\qquad\mathrm{(D)}\ 4\qquad\mathrm{(E)}\ 8$

 \textbf{Answer: }C



\textbf{Solution}

A cube with volume $1$ has a side of length $\sqrt[3]{1}=1$ and thus a surface area of $6 \cdot 1^2=6$. 

 A cube whose surface area is $6\cdot2=12$ has a side of length $\sqrt{\frac{12}{6}}=\sqrt{2}$ and a volume of $(\sqrt{2})^3=2\sqrt{2}\Rightarrow\mathrm{(C)}$.

 
Alternatively, we can use the fact that the surface area of a cube is directly proportional to the square of its side length. Therefore, if the surface area of a cube increases by a factor of $2$, its side length must increase by a factor of $\sqrt{2}$, meaning the new side length of the cube is $1 * \sqrt{2} = \sqrt{2}$. So, its volume is $({\sqrt{2}})^3 = 2\sqrt{2} \Rightarrow\mathrm{(C)}$.

 


	\item (\textit{AMC 12 12A 2008}) Older television screens have an aspect ratio of $4: 3$. That is, the ratio of the width to the height is $4: 3$. The aspect ratio of many movies is not $4: 3$, so they are sometimes shown on a television screen by "letterboxing" - darkening strips of equal height at the top and bottom of the screen, as shown. Suppose a movie has an aspect ratio of $2: 1$ and is shown on an older television screen with a $27$-inch diagonal. What is the height, in inches, of each darkened strip?

                    \begin{figure}[H]
                        \centering
                        \includegraphics[width=0.5\textwidth]{Problem/AMC_12_12A_2008_9_Problem_Figure_1.png}
                    \end{figure}
                    
$\mathrm{(A)}\ 2\qquad\mathrm{(B)}\ 2.25\qquad\mathrm{(C)}\ 2.5\qquad\mathrm{(D)}\ 2.7\qquad\mathrm{(E)}\ 3$

 \textbf{Answer: }D



\textbf{Solution}

Let the width and height of the screen be $4x$ and $3x$ respectively, and let the width and height of the movie be $2y$ and $y$ respectively. 

 By the \href{https://artofproblemsolving.com/wiki/index.php?title=Pythagorean\_Theorem}{Pythagorean Theorem}, the diagonal is $\sqrt{(3x)^2+(4x)^2}=5x = 27$. So $x=\frac{27}{5}$. 

 Since the movie and the screen have the same width, $2y=4x\Rightarrow y=2x$. 

 Thus, the height of each strip is $\frac{3x-y}{2}=\frac{3x-2x}{2}=\frac{x}{2}=\frac{27}{10}=2.7\Longrightarrow\mathrm{(D)}$.

 


 


 



\textbf{Solution 2}

If the ratio is 4:3, and the diagonal is 27, two sides must have been multiplied by a factor of 27/5 to result in the diagonal being 27 (Since 3,4,5 forms a right triangle), multiplying by 27/5, we get width=108/5, height=81/5. From there, we can set up an equation, as it states that after the letterbox has been applied, the aspect ratio became 2:1. 

 We observe that the aspect ratio only takes away from the screen's height, not width. So, if we call X, the height of one darkening strip, we can write this equation: 
 \[\frac{\frac{108}{5}}{\frac{81}{5} - 2x} = 2\]. 

 Therefore, $x = \frac{27}{10} = 2.7$.

 
Thus, the answer is $\text{\textbf{(D)}}$.

 


	\item (\textit{AMC 12 12A 2008}) Doug can paint a room in $5$ hours. Dave can paint the same room in $7$ hours. Doug and Dave paint the room together and take a one-hour break for lunch. Let $t$ be the total time, in hours, required for them to complete the job working together, including lunch. Which of the following equations is satisfied by $t$?


 $\mathrm{(A)}\ \left(\frac{1}{5}+\frac{1}{7}\right)\left(t+1\right)=1\qquad\mathrm{(B)}\ \left(\frac{1}{5}+\frac{1}{7}\right)t+1=1\qquad\mathrm{(C)}\ \left(\frac{1}{5}+\frac{1}{7}\right)t=1\\\mathrm{(D)}\ \left(\frac{1}{5}+\frac{1}{7}\right)\left(t-1\right)=1\qquad\mathrm{(E)}\ \left(5+7\right)t=1$

 \textbf{Answer: }D



\textbf{Solution 1}

Doug can paint $\frac{1}{5}$ of a room per hour, Dave can paint $\frac{1}{7}$ of a room per hour, and the time they spend working together is $t-1$.

 Since rate multiplied by time gives output, $\left(\frac{1}{5}+\frac{1}{7}\right)\left(t-1\right)=1 \Rightarrow \mathrm{(D)}$




\textbf{Solution 2}

If one person does a job in $a$ hours and another person does a job in $b$ hours, the time it takes to do the job together is $\frac{ab}{a+b}$ hours.

 Since Doug paints a room in 5 hours and Dave paints a room in 7 hours, they both paint in $\frac{5\times7}{5+7} = \frac{35}{12}$ hours. They also take 1 hour for lunch, so the total time $t = \frac{35}{12} + 1$ hours.

 Looking at the answer choices, $(D)$ is the only one satisfied by $t = \frac{35}{12} + 1$.

 


	\item (\textit{AMC 12 12A 2008}) Three cubes are each formed from the pattern shown. They are then stacked on a table one on top of another so that the $13$ visible numbers have the greatest possible sum. What is that sum?


 
                    \begin{figure}[H]
                        \centering
                        \includegraphics[width=0.5\textwidth]{Problem/AMC_12_12A_2008_11_Problem_Figure_1.png}
                    \end{figure}
                    
$\mathrm{(A)}\ 154\qquad\mathrm{(B)}\ 159\qquad\mathrm{(C)}\ 164\qquad\mathrm{(D)}\ 167\qquad\mathrm{(E)}\ 189$

 \textbf{Answer: }C



\textbf{Solution}

To maximize the sum of the $13$ faces that are showing, we can minimize the sum of the numbers of the $5$ faces that are not showing. 

 The bottom $2$ cubes each have a pair of opposite faces that are covered up. When the cube is folded, $(1,32)$; $(2,16)$; and $(4,8)$ are opposite pairs. Clearly $4+8=12$ has the smallest sum. 

 The top cube has 1 number that is not showing. The smallest number on a face is $1$. 

 So, the minimum sum of the $5$ unexposed faces is $2\cdot12+1=25$.  Since the sum of the numbers on all the cubes is $3(32+16+8+4+2+1)=189$, the maximum possible sum of $13$ visible numbers is $189-25=164 \Rightarrow C$. 

 



\textbf{Solution 2}

Conversely, maximize the sum. Two cubes have 4 exposed faces. Since $32>16+8+4+2$, 32 must be on the side. There are two distinct (asymmetrical) configurations with 32 on the side, but $(32,16,2,1)$ is the greatest at 51. There are 2 such cubes, so 51*2. The top cube has one unexposed face, so use 1 as the unexposed face. $2(51)+32+16+8+4+2=164 \implies \boxed{\textbf{(C)} 164}$




\textbf{Solution 3}

Two cubes have 4 exposed faces in the same configuration (a ring of 4 squares around the cube). Since the numbers are the same on each cube, they will have the same maximal sum, which is at least 32+16 = 48.The third cube has 5 visible faces. Obviously, the non-visible face on this cube will be the smallest, which is 1. So, this third cube contributes 32+16+8+4+2 = 62. The sum is then at least 62 + 48 + 48. The sum must be even because 62 is even and there are two cubes with the same sum. The only answer choice left is $\boxed{\textbf{(C)} 164}$



	\item (\textit{AMC 12 12A 2008}) A \href{https://artofproblemsolving.com/wiki/index.php?title=Function}{function} $f$ has \href{https://artofproblemsolving.com/wiki/index.php?title=Domain}{domain} $[0,2]$ and \href{https://artofproblemsolving.com/wiki/index.php?title=Range}{range} $[0,1]$. (The notation $[a,b]$ denotes $\{x:a \le x \le b \}$.) What are the domain and range, respectively, of the function $g$ defined by $g(x)=1-f(x+1)$?


 $\mathrm{(A)}\ [-1,1],[-1,0]\qquad\mathrm{(B)}\ [-1,1],[0,1]\qquad\textbf{(C)}\ [0,2],[-1,0]\qquad\mathrm{(D)}\ [1,3],[-1,0]\qquad\mathrm{(E)}\ [1,3],[0,1]$

 \textbf{Answer: }B



\textbf{Solution 1}

$g(x)$ is defined if $f(x + 1)$ is defined. Thus the domain is all $x| x + 1 \in [0,2] \rightarrow x \in [ - 1,1]$. 

 Since $f(x + 1) \in [0,1]$, $- f(x + 1) \in [ - 1,0]$. Thus $g(x) = 1 - f(x + 1) \in [0,1]$ is the range of $g(x)$. 

 Thus the answer is $[- 1,1],[0,1] \longrightarrow \boxed{B}$.

 



\textbf{Solution 2}

Look at the vertical and horizontal translations, and we rewrite the function to a more recognizable $-f(x+1) + 1$ to help us visualize.

 Horizontal: There is one horizontal shift one unit to the left from the $(x+1)$ component making the domain $[-1, 1]$

 Vertical: There is one vertical mirror from the $-f$ causing the range to become $[-1, 0]$ and then a vertical shift one unit upward from the $+ 1$ causing the range to become $[0, 1]$.

 This generates the answer of $\mathrm{(B)}$.

 


	\item (\textit{AMC 12 12A 2008}) Points $A$ and $B$ lie on a circle centered at $O$, and $\angle AOB = 60^\circ$. A second circle is internally tangent to the first and tangent to both $\overline{OA}$ and $\overline{OB}$. What is the ratio of the area of the smaller circle to that of the larger circle?


 $\mathrm{(A)}\ \frac{1}{16}\qquad\mathrm{(B)}\ \frac{1}{9}\qquad\mathrm{(C)}\ \frac{1}{8}\qquad\mathrm{(D)}\ \frac{1}{6}\qquad\mathrm{(E)}\ \frac{1}{4}$

 \textbf{Answer: }B



\textbf{Solution 1}


                    \begin{figure}[H]
                        \centering
                        \includegraphics[width=0.5\textwidth]{Solution/AMC_12_12A_2008_13_Solution_Figure_1.png}
                    \end{figure}
                    Let $P$ be the center of the small circle with radius $r$, and let $Q$ be the point where the small circle is \href{https://artofproblemsolving.com/wiki/index.php?title=Tangent}{tangent} to $OA$. Also, let $C$ be the point where the small circle is tangent to the big circle with radius $R$.

 Then $PQO$ is a right triangle. Angle $POQ$ is $30$ degrees because line $OP$ bisects angle $AOB$ (this can be proved by dropping a perpendicular line from $P$ to line $OB$, letting their intersection be point $S$, and proving triangles $PQO$ and $PSO$ congruent), meaning that $PQO$ is a $30-60-90$ triangle. Therefore, $OP=2PQ$.

 Since $OP=OC-PC=OC-r=R-r$, we have $R-r=2PQ$, or $R-r=2r$, or $\frac{1}{3}=\frac{r}{R}$.

 Ratio of areas of circles is ratio of radii squared, so the answer is $\left(\frac{1}{3}\right)^2 = \frac{1}{9} \Rightarrow \boxed{B}$




\textbf{Solution 2}

Like in Solution 1, let $P$ be the center of the small circle with radius $r$, and let $Q$ be the point where the small circle is \href{https://artofproblemsolving.com/wiki/index.php?title=Tangent}{tangent} to $OA$.

 Let $N$ be the tangency point of the two circles. As shown in Solution 1, $POQ = 30$ degrees, so angle $NOA$ is also $30$ degrees. Let the line tangent to the two circles at $N$ intersect $OA$ and $OB$ at points $C$ and $D$, respectively. Since line $CD$ is tangent to circles $O$ and $P$, it must be perpendicular to $ON$, meaning that angle $ONC$ must be $90$ degrees. Because angle $NOA$ is $30$ degrees, angle $DCO$ is $180-30-90 = 60$ degrees. Angle $DOC$ is also $60$ degrees, so triangle $DOC$ is equilateral.

 Note that an equilateral triangle's incenter is also its centroid. This means the center of the inscribed circle is also the centroid. From properties of median lengths, the radius of the large circle is 3 times the radius of the small circle. 

 Then the ratio of areas will be $\frac{1}{3}$ squared, or $\frac{1}{9}\Rightarrow \boxed{\text{B}}$.

 


	\item (\textit{AMC 12 12A 2008}) What is the area of the region defined by the \href{https://artofproblemsolving.com/wiki/index.php?title=Inequality}{inequality} $|3x-18|+|2y+7|\le3$?


 $\mathrm{(A)}\ 3\qquad\mathrm{(B)}\ \frac {7}{2}\qquad\mathrm{(C)}\ 4\qquad\mathrm{(D)}\ \frac{9}{2}\qquad\mathrm{(E)}\ 5$

 \textbf{Answer: }A



\textbf{Solution 1}

We note that the area of a shape is invariant under translation, meaning that we can translate the inequality left $6$ units and up $\frac{7}{2}$ units without changing the area of the shape defined by the inequality. Doing this, we get a new inequality that is easier to work with:

 
 \[|3x| + |2y|\leq 3\]

 This forms a \href{https://artofproblemsolving.com/wiki/index.php?title=Rhombus}{diamond} centered at the \href{https://artofproblemsolving.com/wiki/index.php?title=Origin}{origin} with vertices at $(\pm 1, 0), (0, \pm 1.5)$. To calculate the area of this diamond, we can split it into four \href{https://artofproblemsolving.com/wiki/index.php?title=Congruence\_(geometry)}{congruent} right triangles. The area of each right triangle is $\frac{1\cdot1.5}{2}=\frac{1.5}{2}$. Multiplying this by $4$ gives us the area of the region defined by the inequality to be $4\cdot \frac{1.5}{2}=\boxed{\textbf{(A) }3}$.

 



\textbf{Solution 2}

Similar to the above solution, but instead we solve for the inequality.

 In the first condition, $3x - 18 \geq 0$, $x \geq 6$, and $2y + 7 \geq 0$, $y \geq -\frac{7}{2}$.Then, the inequality becomes:
 \[3x - 18 + 2y + 7 \leq 3\]
 \[3x + 2y - 11 \leq 3\]
 \[3x + 2y \leq 14\]
 \[y \leq -\frac{3}{2}x + 7\]Along with the constraint that $x \geq 6$ and $y \geq -\frac{7}{2}$, we have a segment of the line $y = -\frac{3}{2}x + 7$ with endpoints $(6, -2)$ and $(7, -\frac{7}{2})$.

 In the second condition, $3x - 18 \geq 0$, $x \geq 6$, and $2y + 7 < 0$, $y < -\frac{7}{2}$.Then, the inequality becomes:
 \[3x - 18 - 2y - 7 \le 3\]
 \[3x - 2y - 25 \le 3\]
 \[3x - 2y \le 28\]
 \[-2y \le -3x + 28\]
 \[y \geq \frac{3}{2}x - 14\]Along with the constraint that $x \geq 6$ and $y < -\frac{7}{2}$, we have a segment of the line $y = \frac{3}{2}x - 14$ with endpoints $(6, -5)$ and $(7, -\frac{7}{2})$.

 In the third condition, $3x - 18 < 0$, $x < 6$, and $2y + 7 \geq 0$, $y \geq -\frac{7}{2}$.Then, the inequality becomes:
 \[-3x + 18 + 2y + 7 \leq 3\]
 \[-3x + 2y + 25 \leq 3\]
 \[-3x + 2y \leq -22\]
 \[2y \leq 3x - 22\]
 \[y \leq \frac{3}{2}x - 11\]Along with the constraint that $x < 6$ and $y \geq -\frac{7}{2}$, we have a segment of the line $y = \frac{3}{2}x - 11$ with endpoints $(5, -\frac{7}{2})$ and $(6, -5)$.

 We may of course still solve for the fourth segment, but it is obvious that the fourth segment would just be another line connecting the gap. Therefore, we see that the shape that is formed by this inequality is a rhombus. Since we know the length of the diagonals are $2$ and $3$ (by the endpoints of the previous solutions), we know that the area of the region is 
 \[\frac{1}{2} \cdot 2 \cdot 3 = 3\]Hence, the answer is $\boxed{A}$.

 


	\item (\textit{AMC 12 12A 2008}) Let $k={2008}^{2}+{2}^{2008}$. What is the units digit of $k^2+2^k$?


 $\mathrm{(A)}\ 0\qquad\mathrm{(B)}\ 2\qquad\mathrm{(C)}\ 4\qquad\mathrm{(D)}\ 6\qquad\mathrm{(E)}\ 8$

 \textbf{Answer: }D



\textbf{Solution}

$k \equiv 2008^2 + 2^{2008} \equiv 8^2 + 2^4 \equiv 4+6 \equiv 0 \pmod{10}$. 

 So, $k^2 \equiv 0 \pmod{10}$. Since $k = 2008^2+2^{2008}$   is a multiple of four and the units digit of powers of two repeat in cycles of four, $2^k \equiv 2^4 \equiv 6 \pmod{10}$. 

 Therefore, $k^2+2^k \equiv 0+6 \equiv 6 \pmod{10}$. So the units digit is $6 \Rightarrow \boxed{D}$.

 \textbf{Note} Another way to get $k \equiv 0 \pmod{10}$ is to find the cycles of the last digit.

 For $2008^2$, we need only be concerned with the last digit $8$ since the other digits do not affect the last digit. Since $8^{2} = 64$, the last digit of $2008^2$ is $4$.

 For $2^{2008}$, note that the last digit cycles through the pattern ${2, 4, 8, 6}$. (You can try to see this by calculating the first powers of $2$.)

 Since $2008$ is a multiple of $4$, the last digit of $2^{2008}$ is evidently $6.$

 Continue as follows (more detail in Solution 2)

 



\textbf{Solution 2}

I am going to share another approach to this problem.

 A units digit $k$ for an integer $n$ implies $n \equiv k \pmod{10}$

 Let us take this step by step. First, we consider $k^2.$

 Note that $k^2 = \left(2008^2 + 2^{2008}\right)^2 = 2008^4 + 2 \cdot 2008^2 \cdot 2^{5016}.$ Now we calculate $k^2 \pmod{10}$

 Before continuing, though, we must take note of the following:

 
 \begin{align*} 2^1 &\equiv 2 \pmod {10} \\ 2^2 &\equiv 4 \pmod {10} \\ 2^3 &\equiv 8 \pmod {10} \\ 2^4 &\equiv 6 \pmod {10} \end{align*}


 Now, we continue with the calculation.

 
 \begin{align*} 2008^4 + 2 \cdot 2008^2 \cdot 2^{5016} &\equiv 8^4 + 2 \cdot 8^2 \cdot 2^{2008} + 2^{5016} \pmod{10} \\ &\equiv 8^4 + 2^1 \cdot 2^6 \cdot 2^{2008} + 2^{5016} \pmod{10} \\ &\equiv 8^4 + 2^{1+6+2008} + 2^{5016} \pmod{10} \\ &\equiv 6 + 2^{2015} + 6 \pmod{10} \\ &\equiv 6 + 2^{3} + 6 \pmod{10} \\ &\equiv 6 + 8 + 6 \pmod{10} \\ &\equiv 20 \pmod{10} \\ &\equiv 0 \pmod{10} \\ \end{align*}


 We do the same with $2^k.$ However, we just need to find $k \pmod 4$ in order to do this calculation since we have the table of $2^k \pmod 10.$

 \begin{align*} 2008^2 + 2^{2008} &\equiv 8^2 \pmod{4} \\ &\equiv 64 \pmod 4\\ &\equiv 0 \pmod 4 \end{align*}


 This implies that 
 \begin{align*} 2^k &\equiv 2^{4} \pmod{10} \\ &\equiv 6 \pmod{10} \end{align*}


 Thus, 
 \begin{align*} k^2 + 2^k &\equiv 6+0 \pmod{10}\\ &\equiv \boxed{\textbf{(D) }6} \pmod{10} \end{align*}





\textbf{Solution 3 (Dangerously Quick)}

By unit digit arithmetic, the unit digit of $k^2+2^k$ need to be either $4$ or $6.$ Hence, we can guess one of them for a $50$\% of getting it right.This should only take 20 seconds or less. ~peelybonehead

 Note: This solution is not recommended - this is effectively a truncated version of Solution 1 - do this only if you have very little time left.

 


	\item (\textit{AMC 12 12A 2008}) The numbers $\log(a^3b^7)$, $\log(a^5b^{12})$, and $\log(a^8b^{15})$ are the first three terms of an arithmetic sequence, and the $12^\text{th}$ term of the sequence is $\log(b^n)$. What is $n$?


 $\mathrm{(A)}\ 40\qquad\mathrm{(B)}\ 56\qquad\mathrm{(C)}\ 76\qquad\mathrm{(D)}\ 112\qquad\mathrm{(E)}\ 143$

 \textbf{Answer: }D



\textbf{Solution 1}

Let $A = \log(a)$ and $B = \log(b)$. 

 The first three terms of the arithmetic sequence are $3A + 7B$, $5A + 12B$, and $8A + 15B$, and the $12^\text{th}$ term is $nB$. 

 Thus, $2(5A + 12B) = (3A + 7B) + (8A + 15B) \Rightarrow A = 2B$. 

 Since the first three terms in the sequence are $13B$, $22B$, and $31B$, the $k$th term is $(9k + 4)B$.

 Thus the $12^\text{th}$ term is $(9\cdot12 + 4)B = 112B = nB \Rightarrow n = 112\Rightarrow \boxed{D}$.

 



\textbf{Solution 2}

If $\log(a^3b^7)$, $\log(a^5b^{12})$, and $\log(a^8b^{15})$ are in \href{https://artofproblemsolving.com/wiki/index.php?title=Arithmetic\_progression}{arithmetic progression}, then $a^3b^7$, $a^5b^{12}$, and $a^8b^{15}$ are in \href{https://artofproblemsolving.com/wiki/index.php?title=Geometric\_progression}{geometric progression}. Therefore,

 
 \[a^2b^5=a^3b^3 \Rightarrow a=b^2\]

 Therefore, $a^3b^7=b^{13}$, $a^5b^{12}=b^{22}$, therefore the 12th term in the sequence is $b^{13+9*11}=b^{112} \Rightarrow \boxed{D}$




\textbf{Solution 3}

$\ \text{If a, b, and c are in a arithmetic progression then } b = \frac{a+c}{2} \text{ which means}$$\ \log(a^5b^{12}) = \frac{\log(a^3b^7) + \log(a^8b^{15})}{2}= \frac{\log(a^{11}b^{22})}{2} \text{ therefore}$$\ 2\log(a^5b^{12}) = \log(a^{10}b^{24}) = \log(a^{11}b^{22}) \Rightarrow a=b^2$$\ \text{This means that the Kth term of the series would be } \log(b^{13+9(k-1)})$$\ \text{The 12th term would be } \log(b^{112})  \Rightarrow n=112  \Rightarrow D$




\textbf{Solution 4 (mimimal manipulation)}

Given the first three terms form an arithmetic progression, we have:
 \[a = \log(a^3b^7)\]
 \[a+d = \log(a^5b^{12})\]
 \[a+2d = \log(a^8b^{15}).\]Subtracting the first equation from the second and the third from the second, respectively, gives us these two expressions for $d$:
 \[d = \log \left( \frac{a^5b^{12}}{a^3b^7} \right) = \log a^2b^5\]
 \[d = \log \left( \frac{a^8b^{15}}{a^5b^{12}} \right) = \log a^3b^3.\]The desired $12$th term in the sequence is $a+11d$, so we can substitute our values for $a$ and $d$ (using either one of our two expressions for $d$):
 \[a+11d = \log a^3b^7 + 11\log(a^2b^5)\]
 \[= \log a^3b^7 + \log(a^{22}b^{55})\]
 \[= \log a^{25}b^{62}.\]The answer must be expressed as $\log(b^n)$, however. We're in luck: the two different yet equal expressions for $d$ allow us to express $a$ and $b$ in terms of each other:
 \[\log a^2b^5 = \log a^3b^3\]
 \[a^2b^5 = a^3b^3\]
 \[a=b^2.\]Plugging in $a=b^2$, we have:
 \[a+11d = \log b^{50}b^{62}\]
 \[= \log b^{112} \Rightarrow \boxed{D}.\]~ Jingwei325 $\smiley$



	\item (\textit{AMC 12 12A 2008}) Let $a_1,a_2,\ldots$ be a sequence determined by the rule $a_n=a_{n-1}/2$ if $a_{n-1}$ is even and $a_n=3a_{n-1}+1$ if $a_{n-1}$ is odd. For how many positive integers $a_1 \le 2008$ is it true that $a_1$ is less than each of $a_2$, $a_3$, and $a_4$? 


 $\mathrm{(A)}\ 250\qquad\mathrm{(B)}\ 251\qquad\mathrm{(C)}\ 501\qquad\mathrm{(D)}\ 502\qquad\mathrm{(E)} 1004$

 \textbf{Answer: }D



\textbf{Solution 1}

All positive integers can be expressed as $4n$, $4n+1$, $4n+2$, or $4n+3$, where $n$ is a nonnegative integer. 

 
            \begin{itemize}
        \item If $a_1=4n$, then $a_2=\frac{4n}{2}=2n<a_1$.
            \end{itemize}
            
            \begin{itemize}
        \item If $a_1=4n+1$, then $a_2=3(4n+1)+1=12n+4$, $a_3=\frac{12n+4}{2}=6n+2$, and $a_4=\frac{6n+2}{2}=3n+1<a_1$.
            \end{itemize}
            
            \begin{itemize}
        \item If $a_1=4n+2$, then $a_2=2n+1<a_1$.
            \end{itemize}
            
            \begin{itemize}
        \item If $a_1=4n+3$, then $a_2=3(4n+3)+1=12n+10$, $a_3=\frac{12n+10}{2}=6n+5$, and $a_4=3(6n+5)+1=18n+16$.
            \end{itemize}
            Since $12n+10, 6n+5, 18n+16 > 4n+3$, every positive integer $a_1=4n+3$ will satisfy $a_1<a_2,a_3,a_4$. 

 Since one fourth of the positive integers $a_1 \le 2008$ can be expressed as $4n+3$, where $n$ is a nonnegative integer, the answer is $\frac{1}{4}\cdot 2008 = 502 \Rightarrow D$.

 



\textbf{Solution 2}

After checking the first few $a_n$ such as $1$, $2$ through $7$, we can see that the only $a_1$ that satisfy the conditions are odd numbers that when tripled and added 1 to, are double an odd number. For example, for $a_n=3$, we notice the sequence yields $10$, $5$, and $16$, a valid sequence.

 So we can set up an equation, $3x + 1 = 2(2k - 1)$ where x is equal to $a_1$. Rearranging the equation yields $(3x + 3)/4 = k$. Experimenting yields that every 4th $x$ after 3 creates an integer, and thus satisfies the sequence condition. So the number of valid solutions is equal to $2008/4 = 502 \Rightarrow D$.

 



\textbf{Solution 3}

First, we can easily see that $a_1$ cannot be even since for all positive integers $x$, we have $x>\dfrac{x}2$. That means $a_1$ must be odd, so $a_2=3a_1+1$. That means $a_2$ must be even, so $a_3=\dfrac{a_2}2=\dfrac{3a_1+1}2$. Now, let's see if $a_2$ and $a_3$ are greater than $a_1$. We have $a_2=3a_1+1$, and we can easily see that $3x+1>x$ for all positive integers $x$, so for all odd positive integers $a_1$, we have $a_2>a_1$. Now, let's see if $a_3>a_1$. If $a_3>a_1$, then $\dfrac{3a_1+1}2>a_1$, so $3a_1+1>2a_1$, so $a_1>-1$, and that is true for all odd positive integers $a_1$, so we would have $a_3>a_1$ for all odd positive integers $a_1$. Now, we have two cases!

 \textbf{Case 1: $\mathbf{a_3}$ is odd} \noindent That means $a_4=3a_3+1$, so $a_4>a_3$ must be true. Since we also know that $a_3>a_1$ is true, we must have $a_4>a_1$. Therefore, this case always works! Now, let's see when would $a_3$ be odd. If $a_3$ is odd, then $a_2$ should contain exactly $1$ two. That means we must have $a_2$ be $2,6,10,14,18,22,26,\cdots$. For there to exist an $a_1$, $a_2$ must be one more than a multiple of $3$, so we must have $a_2=10,22,34,46,\ldots$ which corresponds to $a_1=3,7,11,15,\ldots,2007$, so this case gives $502$ possible values of $a_1$.

 \textbf{Case 1: $\mathbf{a_3}$ is even} \noindent That means $a_4=\dfrac{a_3}2$, so 
 \[a_4=\dfrac{\left(\dfrac{3a_1+1}2\right)}2=\dfrac{3a_1+1}4.\]If $a_4>a_1$, then $\dfrac{3a_1+1}4>a_1$, so $3a_1+1>4a_1$, so $1>a_1$. But, since $a_1$ must be a positive integer, there cannot exist any possible $a_1$ such that $1>a_1$, so this case will not give any possible values of $a_1$.\noindent\textbf{Conclusion}

 \noindent Totaling up our two cases, we see that $502+0=\boxed{\mathrm{(D)}\ 502}$ possible positive integers $a_1$ exist.

 


	\item (\textit{AMC 12 12A 2008}) A triangle $\triangle ABC$ with sides $5$, $6$, $7$ is placed in the three-dimensional plane with one vertex on the positive $x$ axis, one on the positive $y$ axis, and one on the positive $z$ axis. Let $O$ be the origin. What is the volume of $OABC$?  


 $\mathrm{(A)}\ \sqrt{85}\qquad\mathrm{(B)}\ \sqrt{90}\qquad\mathrm{(C)}\ \sqrt{95}\qquad\mathrm{(D)}\ 10\qquad\mathrm{(E)}\ \sqrt{105}$

 \textbf{Answer: }C



\textbf{Solution}


                    \begin{figure}[H]
                        \centering
                        \includegraphics[width=0.5\textwidth]{Solution/AMC_12_12A_2008_18_Solution_Figure_1.png}
                    \end{figure}
                    

 Without loss of generality, let $A$ be on the $x$ axis, $B$ be on the $y$ axis, and $C$ be on the $z$ axis, and let $AB, BC, CA$ have respective lengths of 5, 6, and 7.  Let $a,b,c$ denote the lengths of segments $OA,OB,OC,$ respectively.  Then by the \href{https://artofproblemsolving.com/wiki/index.php?title=Pythagorean\_Theorem}{Pythagorean Theorem},
 \begin{align*} a^2+b^2 &=5^2 , \\  b^2+c^2&=6^2, \\ c^2+a^2 &=7^2 , \end{align*}
so $a^2 = (5^2+7^2-6^2)/2 = 19$; similarly, $b^2 = 6$ and $c^2 = 30$. Since $OA$, $OB$, and $OC$ are mutually perpendicular, the tetrahedron's volume is 
 \[abc/6\] because we can consider the tetrahedron to be a right triangular pyramid. 
 \[abc/6 = \sqrt{a^2b^2c^2}/6 = \frac{\sqrt{19 \cdot 6 \cdot 30}}{6} = \sqrt{95},\]which is answer choice $\boxed{\text{C}}$.

 


	\item (\textit{AMC 12 12A 2008}) In the expansion of

 \[\left(1 + x + x^2 + \cdots + x^{27}\right)\left(1 + x + x^2 + \cdots + x^{14}\right)^2,\]
what is the \href{https://artofproblemsolving.com/wiki/index.php?title=Coefficient}{coefficient} of $x^{28}$?


 $\mathrm{(A)}\ 195\qquad\mathrm{(B)}\ 196\qquad\mathrm{(C)}\ 224\qquad\mathrm{(D)}\ 378\qquad\mathrm{(E)}\ 405$

 \textbf{Answer: }C



\textbf{Solution 1 (easiest)}

Let $A = \left(1 + x + x^2 + \cdots + x^{14}\right)$ and $B = \left(1 + x + x^2 + \cdots + x^{27}\right)$. We are expanding $A \cdot A \cdot B$. 

 Since there are $15$ terms in $A$, there are $15^2 = 225$ ways to choose one term from each $A$. The product of the selected terms is $x^n$ for some integer $n$ between $0$ and $28$ inclusive. For each $n \neq 0$, there is one and only one $x^{28 - n}$ in $B$. For example, if I choose $x^2$ from $A$ , then there is exactly one power of $x$ in $B$ that I can choose; in this case, it would be $x^{24}$. Since there is only one way to choose one term from each $A$ to get a product of $x^0$, there are $225 - 1 = 224$ ways to choose one term from each $A$ and one term from $B$ to get a product of $x^{28}$. Thus the coefficient of the $x^{28}$ term is $224 \Rightarrow \boxed{C}$.

 



\textbf{Solution 2}

Let $P(x) = \left(1 + x + x^2 + \cdots + x^{14}\right)^2 = a_0 + a_1x + a_2x^2 + \cdots + a_{28}x^{28}$. Then the $x^{28}$ term from the product in question $\left(1 + x + x^2 + \cdots + x^{27}\right)(a_0 + a_1x + a_2x^2 + \cdots + a_{28}x^{28})$ is

 $1a_{28}x^{28} + xa_{27}x^{27} + x^2a_{26}x^{26} + \cdots + x^{27}a_1x = \left(a_1 + a_2 + \cdots a_{28}\right)x^{28}$

 So we are trying to find the sum of the coefficients of $P(x)$ minus $a_0$. Since the constant term $a_0$ in $P(x)$ (when expanded) is $1$, and the sum of the coefficients of $P(x)$ is $P(1)$, we find the answer to be$P(1) - a_0 = \left(1 + 1 + 1^2 + \cdots 1^{14}\right)^2 - 1 = 15^2 - 1 = 224 \Rightarrow \boxed{C}$.

 



\textbf{Solution 3}

We expand $(1 + x + x^2 + x^3 + \cdots + x^{14})^2$ to $(1 + x + x^2 + x^3 + \cdots + x^{14}) * (1 + x + x^2 + x^3 + \cdots + x^{14})$ and use FOIL to multiply. It expands out to:

 $1 + x + x^2 + x^3 + x^4 + \cdots + x^{14} +$
$\qquad x + x^2 + x^3 + x^4 + \cdots + x^{14} + x^{15} +$
$\qquad \qquad x^2 + x^3 + x^4 + \cdots + x^{14} + x^{15} + x^{16} + \cdots$

 It becomes apparent that 

 $(1 + x + x^2 + x^3 + \cdots + x^{14})^2 = 1 + 2x + 3x^2 + 4x^3 + \cdots + 15x^{14} + 14x^{15} + 13x^{16} + \cdots + x^{28}$.

 Now we have to find the coefficient of $x^{28}$ in the product: 

 $(1 + 2x + 3x^2 + 4x^3 + \cdots + 15x^{14} + 14x^{15} + 13x^{16} + \cdots + x^{28}) \cdot (1 + x + x^2 + x^3 + \cdots + x^{27})$.

 We quickly see that the we get $x^{28}$ terms from $x^{27} \cdot 2x$, $x^{26} \cdot 3x^2$, $x^{25} \cdot 4x^3$, ... $15x^{14} \cdot x^{14}$, ... $x^{28} \cdot 1$. The coefficient of $x^{28}$ is just the sum of the coefficients of all these terms. $1 + 2 + 3 + 4 + \cdots + 15 + 14 + 13 + \cdots + 4 + 3 + 2 = 224$, so the answer is $\boxed{C}$.

 



\textbf{Solution 4}

Rewrite the product as $\frac{(x^{28} - 1)(x^{15} - 1)(x^{15} - 1)}{(x - 1)^3}$. It is known that 

 
 \[\frac{1}{(1 - x)^n} = \binom{n - 1}{n - 1} + \binom{n}{n - 1}x + \binom{n + 1}{n - 1}x^2 + \binom{n + 2}{n - 1}x^3 + \cdots + \binom{n - 1 + k}{n - 1}x^k + \cdots .\]

 Thus, our product becomes

 
 \[-\left( x^{28} - 1 \right) \left( x^{15} - 1 \right) \left( x^{15} - 1 \right) \left( \binom{2}{2} + \binom{3}{2}x + \binom{4}{2}x^2 + \cdots \right).\]

 \[= -\left( x^{28} - 1 \right) \left( x^{15} - 1 \right) \left( x^{15} - 1 \right) \left( 1 + 3x + 6x^2 + \cdots \right).\]

 We determine the $x^{28}$ coefficient by doing casework on the first three terms in our product. We can obtain an $x^{28}$ term by choosing $x^{28}$ in the first term, $-1$ in the second and third terms, and $1$ in the fourth term. We can get two $x^{28}$ terms by choosing $x^{15}$ in either the second or third term, $-1$ in the first term, $-1$ in the second or third term from which $x^{15}$ has not been chosen, and the $\binom{15}{2}x^{13}$ in the fourth term. We get $\binom{15}{2} * 2 = 210$ $x^{28}$ terms this way. (We multiply by $2$ because the $x^{15}$ term could have been chosen from the second term or the third term). Lastly, we can get an $x^{28}$ term by choosing $-1$ in the first three terms and a $\binom{30}{2}x^{28}$ from the fourth term. We have a total of $1 + 210 - 435 = -224$ for the $x^{28}$ coefficient, but we recall that we have a negative sign in front of our product, so we obtain an answer of $224 \Rightarrow \boxed{(C)}$.

 



\textbf{Solution 5}

We can use the Stars and Bars Formula to create an equation$a + b + c = 28$, with restrictions $a <= 27, b<=14, and c<= 14$. We also note that any variable can have a power of 0. Using the formula (n + k - 1 choose k -1), we find that the total possibilities of getting an exponent of 27 is...$(28 + 2 choose 2) - [(2 choose 2) + (28 - 15 + 2 choose 2) + (28 - 15 + 2 choose 2)].$We stop since it is impossible for both variables to exceed the limit.Therefore, simplifying out equation we get $(30 choose 2) - [1 + 2(15 choose 2)] = 435 - 211 = 224$, correcponding to choice $C$!

 


	\item (\textit{AMC 12 12A 2008}) Triangle $ABC$ has $AC=3$, $BC=4$, and $AB=5$. Point $D$ is on $\overline{AB}$, and $\overline{CD}$ bisects the right angle. The inscribed circles of $\triangle ADC$ and $\triangle BCD$ have radii $r_a$ and $r_b$, respectively. What is $r_a/r_b$?


 $\mathrm{(A)}\ \frac{1}{28}\left(10-\sqrt{2}\right)\qquad\mathrm{(B)}\ \frac{3}{56}\left(10-\sqrt{2}\right)\qquad\mathrm{(C)}\ \frac{1}{14}\left(10-\sqrt{2}\right)\qquad\mathrm{(D)}\ \frac{5}{56}\left(10-\sqrt{2}\right)\\\mathrm{(E)}\ \frac{3}{28}\left(10-\sqrt{2}\right)$

 \textbf{Answer: }E



\textbf{Solution 1}


                    \begin{figure}[H]
                        \centering
                        \includegraphics[width=0.5\textwidth]{Solution/AMC_12_12A_2008_20_Solution_Figure_1.png}
                    \end{figure}
                    

 By the \href{https://artofproblemsolving.com/wiki/index.php?title=Angle\_Bisector\_Theorem}{Angle Bisector Theorem}, 
 \[\frac{BD}{4} = \frac{5-BD}{3} \Longrightarrow BD = \frac{20}7\]By \href{https://artofproblemsolving.com/wiki/index.php?title=Law\_of\_Sines}{Law of Sines} on $\triangle BCD$, 
 \[\frac{BD}{\sin 45^{\circ}} = \frac{CD}{\sin \angle B} \Longrightarrow \frac{20/7}{\sqrt{2}/2} = \frac{CD}{3/5} \Longrightarrow CD=\frac{12\sqrt{2}}{7}\]Since the area of a triangle satisfies $[\triangle]=rs$, where $r =$ the \href{https://artofproblemsolving.com/wiki/index.php?title=Inradius}{inradius} and $s =$ the \href{https://artofproblemsolving.com/wiki/index.php?title=Semiperimeter}{semiperimeter}, we have 
 \[\frac{r_A}{r_B} = \frac{[ACD] \cdot s_B}{[BCD] \cdot s_A}\]Since $\triangle ACD$ and $\triangle BCD$ share the \href{https://artofproblemsolving.com/wiki/index.php?title=Altitude}{altitude} (to $\overline{AB}$), their areas are the ratio of their bases, or 
 \[\frac{[ACD]}{[BCD]} = \frac{AD}{BD} = \frac{3}{4}\]The semiperimeters are $s_A = \left(3 + \frac{15}{7} + \frac{12\sqrt{2}}{7}\right)\left/\right.2 = \frac{18+6\sqrt{2}}{7}$ and $s_B = \frac{24+ 6\sqrt{2}}{7}$. Thus,
 \begin{align*} \frac{r_A}{r_B} &= \frac{[ACD] \cdot s_B}{[BCD] \cdot s_A} = \frac{3}{4} \cdot \frac{(24+ 6\sqrt{2})/7}{(18+6\sqrt{2})/7} \\ &= \frac{3(4+\sqrt{2})}{4(3+\sqrt{2})} \cdot \left(\frac{3-\sqrt{2}}{3-\sqrt{2}}\right) = \frac{3}{28}(10-\sqrt{2}) \Rightarrow \mathrm{(E)}\qquad \blacksquare \end{align*}


 -(you could also use stewart's theorem to find the length of cd)

 



\textbf{Solution 2}


                    \begin{figure}[H]
                        \centering
                        \includegraphics[width=0.5\textwidth]{Solution/AMC_12_12A_2008_20_Solution_Figure_2.png}
                    \end{figure}
                    

 We start by finding the length of $AD$ and $BD$ as in solution 1. Using the angle bisector theorem, we see that $AD = \frac{15}{7}$ and $BD = \frac{20}{7}$. Using Stewart's Theorem gives us the equation $5d^2 + \frac{1500}{49} = \frac{240}{7} + \frac{180}{7}$, where $d$ is the length of $CD$. Solving gives us $d = \frac{12\sqrt{2}}{7}$, so $CD = \frac{12\sqrt{2}}{7}$.

 Call the incenters of triangles $ACD$ and $BCD$ $O_a$ and $O_b$ respectively. Since $O_a$ is an incenter, it lies on the angle bisector of $\angle ACD$. Similarly, $O_b$ lies on the angle bisector of $\angle BCD$. Call the point on $CD$ tangent to $O_a$ $M$, and the point tangent to $O_b$ $N$. Since $\triangle CO_aM$ and $\triangle CO_bN$ are right, and $\angle O_aCM = \angle O_bCN$, $\triangle CO_aM \sim \triangle CO_bN$. Then, $\frac{r_a}{r_b} = \frac{CM}{CN}$. 

 We now use common tangents to find the length of $CM$ and $CN$. Let $CM = m$, and the length of the other tangents be $n$ and $p$. Since common tangents are equal, we can write that $m + n = \frac{12\sqrt{2}}{7}$, $n + p = \frac{15}{7}$ and $m + p = 3$. Solving gives us that $CM = m = \frac{6\sqrt{2} + 3}{7}$. Similarly, $CN = \frac{6\sqrt{2} + 4}{7}$. 

 We see now that $\frac{r_a}{r_b} = \frac{\frac{6\sqrt{2} + 3}{7}}{\frac{6\sqrt{2} + 4}{7}} = \frac{6\sqrt{2} + 3}{6\sqrt{2} + 4} = \frac{60-6\sqrt{2}}{56} = \frac{3}{28}(10 - \sqrt{2}) \Rightarrow \boxed{E}$



	\item (\textit{AMC 12 12A 2008}) A permutation $(a_1,a_2,a_3,a_4,a_5)$ of $(1,2,3,4,5)$ is $\textit{heavy-tailed}$ if $a_1 + a_2 < a_4 + a_5$.  What is the number of heavy-tailed permutations?


 $\mathrm{(A)}\ 36\qquad\mathrm{(B)}\ 40\qquad\textbf{(C)}\ 44\qquad\mathrm{(D)}\ 48\qquad\mathrm{(E)}\ 52$

 \textbf{Answer: }D



\textbf{Solution 1 (Symmetry)}

There are $5!=120$ total permutations. 

 For every permutation $(a_1,a_2,a_3,a_4,a_5)$ such that $a_1 + a_2 < a_4 + a_5$, there is exactly one permutation such that $a_1 + a_2 > a_4 + a_5$. Thus it suffices to count the permutations such that $a_1 + a_2 = a_4 + a_5$. 

 $1+4=2+3$, $1+5=2+4$, and $2+5=3+4$ are the only combinations of numbers that can satisfy $a_1 + a_2 = a_4 + a_5$. 

 There are $3$ combinations of numbers, $2$ possibilities of which side of the equation is $a_1+a_2$ and which side is $a_4+a_5$, and $2^2=4$ possibilities for rearranging $a_1,a_2$ and $a_4,a_5$. Thus, there are $3\cdot2\cdot4=24$ permutations such that $a_1 + a_2 = a_4 + a_5$. 

 Thus, the number of <u>heavy-tailed</u> permutations is $\frac{120-24}{2}=48 \Rightarrow D$.

 



\textbf{Solution 2 (Casework)}

We use case work on the value of $a_3$. 

 Case 1: $a_3 = 1$. Since $a_1 + a_2 < a_4 + a_5$, $(a_1, a_2)$ can only be a permutation of $(2, 3)$ or $(2, 4)$. The values of $a_1$ and $a_2$, as well as the values of $a_4$ and $a_5$, are interchangeable, so this case produces a total of $2(2 \cdot 2) = 8$ solutions.

 Case 2: $a_3 = 2$. Similarly, we have $(a_1, a_2)$ is a permutation of $(1, 3)$, $(1, 4)$, or $(1, 5)$, which gives a total of $3(2 \cdot 2) = 12$ solutions.

 Case 3: $a_3 = 3$. $(a_1, a_2)$ is a permutation of $(1, 2)$ or $(1, 4)$, which gives a total of $2(2 \cdot 2) = 8$ solutions.

 Case 4: $a_3 = 4$. $(a_1, a_2)$ is a permutation of $(1, 2)$, $(1, 3)$, or $(2, 3)$, which gives a total of $3(2 \cdot 2) = 12$ solutions. 

 Case 5: $a_3 = 5$. $(a_1, a_2)$ is a permutation of $(1, 2)$ or $(1, 3)$, which gives a total of $2(2 \cdot 2) = 8$ solutions.

 Therefore, our answer is $8 + 12 + 8 + 12 + 8 = 48 \Rightarrow \boxed{D}$.

 -MP8148

 


	\item (\textit{AMC 12 12A 2008}) A round table has radius $4$. Six rectangular place mats are placed on the table. Each place mat has width $1$ and length $x$ as shown. They are positioned so that each mat has two corners on the edge of the table, these two corners being end points of the same side of length $x$. Further, the mats are positioned so that the inner corners each touch an inner corner of an adjacent mat. What is $x$?


 
                    \begin{figure}[H]
                        \centering
                        \includegraphics[width=0.5\textwidth]{Problem/AMC_12_12A_2008_22_Problem_Figure_1.png}
                    \end{figure}
                    
$\mathrm{(A)}\ 2\sqrt{5}-\sqrt{3}\qquad\mathrm{(B)}\ 3\qquad\mathrm{(C)}\ \frac{3\sqrt{7}-\sqrt{3}}{2}\qquad\mathrm{(D)}\ 2\sqrt{3}\qquad\mathrm{(E)}\ \frac{5+2\sqrt{3}}{2}$

 \textbf{Answer: }C



\textbf{Solution 1 (Trigonometry)}

Let one of the mats be $ABCD$, and the center be $O$ as shown: 

 
                    \begin{figure}[H]
                        \centering
                        \includegraphics[width=0.5\textwidth]{Solution/AMC_12_12A_2008_22_Solution_Figure_1.png}
                    \end{figure}
                    

 Since there are $6$ mats, $\Delta BOC$ is \href{https://artofproblemsolving.com/wiki/index.php?title=Equilateral}{equilateral} (the hexagon with side length $x$ is regular). So, $BC=CO=x$. Also, $\angle OCD = \angle OCB + \angle BCD = 60^\circ+90^\circ=150^\circ$. 

 By the \href{https://artofproblemsolving.com/wiki/index.php?title=Law\_of\_Cosines}{Law of Cosines}: $4^2=1^2+x^2-2\cdot1\cdot x \cdot \cos(150^\circ) \Rightarrow x^2 + x\sqrt{3} - 15 = 0 \Rightarrow x = \frac{-\sqrt{3}\pm 3\sqrt{7}}{2}$. 

 Since $x$ must be positive, $x = \frac{3\sqrt{7}-\sqrt{3}}{2} \Rightarrow (C)$.

 



\textbf{Solution 2 (without trigonometry)}

Draw $OD$ and $OC$ as in the diagram. Draw the altitude from $O$ to $DC$ and call the intersection $E$.

 

                    \begin{figure}[H]
                        \centering
                        \includegraphics[width=0.5\textwidth]{Solution/AMC_12_12A_2008_22_Solution_Figure_2.png}
                    \end{figure}
                    

 As proved in the first solution, $\angle OCD = 150^\circ$. That makes $\Delta OCE$ a $30-60-90$ triangle, so $OE = \frac{x}{2}$ and $CE= \frac{x\sqrt 3}{2}$

 Since $\Delta ODE$ is a right triangle, $\left({\frac{x}{2}}\right)^2 + \left({\frac{x\sqrt 3}{2} +1}\right)^2 = 4^2 \Rightarrow  x^2+x\sqrt3-15 = 0$

 Solving for $x$ gives $x =\frac{3\sqrt{7}-\sqrt{3}}{2}\Rightarrow (C)$




\textbf{Solution 3 (simply Pythagorean Theorem)}


                    \begin{figure}[H]
                        \centering
                        \includegraphics[width=0.5\textwidth]{Solution/AMC_12_12A_2008_22_Solution_Figure_3.png}
                    \end{figure}
                    

 By symmetry, $E$ is the midpoint of $DF$ and $OE$ is an extension of $OC$. Thus $\angle OED = 90^\circ$. Since $OD = 4$ and $DE = \frac{1}{2}$, $OE = \sqrt{16-\frac{1}{4}} = \frac{\sqrt{63}}{2} = \frac{3\sqrt{7}}{2}$. Since $\triangle CED$ is $30-60-90$, $CE = \frac{\sqrt{3}}{2}$ (this can also be deduced from Pythagoras on $\triangle CED$).

 Thus $OC = \frac{3\sqrt{7}-\sqrt{3}}{2}$. As previous solutions noted, $\triangle BOC$ is equilateral, and thus the desired length is $x = OC \implies (C)$.

 



\textbf{Solution 4}


                    \begin{figure}[H]
                        \centering
                        \includegraphics[width=0.5\textwidth]{Solution/AMC_12_12A_2008_22_Solution_Figure_4.png}
                    \end{figure}
                    

 Looking at the diagram above, we know that $BE$ is a diameter of circle $O$ due to symmetry. Due to \href{https://artofproblemsolving.com/wiki/index.php?title=Thales\%27\_theorem}{Thales' theorem}, triangle $ABE$ is a right triangle with $A = 90 ^\circ$. $AE$ lies on $AD$ and $GE$ because $BAD$ is also a right angle. To find the length of $DG$, notice that if we draw a line from $F$ to $M$, the midpoint of line $DG$, it creates two $30$ - $60$ - $90$ triangles. Therefore, $MD = MG = \frac{\sqrt{3}x}{2} \Rightarrow DG = \sqrt{3}x$. $AE = 2 + \sqrt{3}x$

 Applying the Pythagorean theorem to triangle $ABE$, we get 
 \[(2+\sqrt{3}x)^2 + x^2 = 8^2 \Rightarrow 4 + 3x^2 + 4\sqrt{3}x + x^2 = 64 \Rightarrow x^2 + \sqrt{3}x - 15 = 0\] Using the quadratic formula to solve, we get 
 \[x = \frac{-\sqrt{3} \pm \sqrt{3 -4(1)(-15)}}{2} = \frac{\pm 3\sqrt{7} - \sqrt{3}}{2}\] $x$ must be positive, therefore 
 \[x = \frac{3\sqrt{7} - \sqrt{3}}{2} \Rightarrow (C)\]




\textbf{Solution 5 (coordinate bashing)}


                    \begin{figure}[H]
                        \centering
                        \includegraphics[width=0.5\textwidth]{Solution/AMC_12_12A_2008_22_Solution_Figure_5.png}
                    \end{figure}
                    

 We will let $O(0,0)$ be the origin. This way the coordinates of $C$ will be $(0,y)$. By $30-60-90$, the coordinates of $D$ will be $\left(-\frac{1}{2}, y + \frac{\sqrt{3}}{2}\right)$. The distance any point with coordinates $(x, y)$ is from the origin is $\sqrt{x^2 + y^2}$. Therefore, the distance $D$ is from the origin is $4$ and $\frac{1}{4} + x^2 + x\sqrt{3} + \frac{3}{4} = x^2 + x\sqrt{3} + 1 = 16$. We get the quadratic equation mentioned in solution 2. Using the quadratic formula, we get that $x = \frac{3\sqrt{7}-\sqrt{3}}{2} \Rightarrow (C)$

 Note: Since $C$ and $D$ are not labeled in the diagram, refer to solution 1 for the location of points $C$ and $D$.

 



\textbf{Solution 6}


                    \begin{figure}[H]
                        \centering
                        \includegraphics[width=0.5\textwidth]{Solution/AMC_12_12A_2008_22_Solution_Figure_6.png}
                    \end{figure}
                    

 Notice that $\overarc{AE}$ is $\frac16$ the circumference of the circle. Therefore, $\overline{AE}$ is the side length of an inscribed hexagon with side length $4$. $\triangle AFE$ is a right triangle with $\overline{AF}=\frac12$. The length of $\overline{EF}=x+\frac{\sqrt{3}}{2}$. Using the \href{https://artofproblemsolving.com/wiki/index.php?title=Pythagorean\_Theorem}{Pythagorean Theorem}, we get

 
$4^2 = \left(\frac{1}{2}\right)^2 + \left(x+\frac{\sqrt{3}}{2}\right)^2$

 
Solving for $x$, we get $x = \frac{3\sqrt{7}-\sqrt{3}}{2}\ \boxed{\text{C}}$




\textbf{Solution 7 (Answer Choices & Estimation)}

The smallest distance between the intersection between two adjacent placemats and the circle is a bit less than $1$. Thus, the answer will be a bit more than $4-1=3.$ Going up, we guess that the next one up is the answer, $\frac{3\sqrt{7}-\sqrt{3}}{2},$ which is about $3.1$. The next one up after that is $2\sqrt{3},$ which is about $3.5$ and seems too big. So we circle $\boxed{\text{(C) }\frac{3\sqrt{7}-\sqrt{3}}{2}}$ and are happy we just cheesed the $AMC$ $10$ $2008$ problem $25.$



	\item (\textit{AMC 12 12A 2008}) The solutions of the equation $z^4+4z^3i-6z^2-4zi-i=0$ are the vertices of a convex polygon in the complex plane. What is the area of the polygon?


 $\mathrm{(A)}\ 2^{\frac{5}{8}}\qquad\mathrm{(B)}\ 2^{\frac{3}{4}}\qquad\mathrm{(C)}\ 2\qquad\mathrm{(D)}\ 2^{\frac{5}{4}}\qquad\mathrm{(E)}\ 2^{\frac{3}{2}}$

 \textbf{Answer: }D



\textbf{Solution 1}

Looking at the coefficients, we are immediately reminded of the binomial expansion of ${\left(x+1\right)}^{4}$.

 Modifying this slightly, we can write the given equation as:
 \[{\left(z+i\right)}^{4}=1+i=2^{\frac{1}{2}}\cos \frac {\pi}{4} + 2^{\frac{1}{2}}i\sin \frac {\pi}{4}\]We can apply a translation of $-i$ and a rotation of $-\frac{\pi}{4}$ (both operations preserve area) to simplify the problem:
 \[z^{4}=2^{\frac{1}{2}}\]

 Because the roots of this equation are created by rotating $\frac{\pi}{2}$ radians successively about the origin, the quadrilateral is a square.

 We know that half the diagonal length of the square is ${\left(2^{\frac{1}{2}}\right)}^{\frac{1}{4}}=2^{\frac{1}{8}}$

 Therefore, the area of the square is$\frac{{\left( 2 \cdot 2^{\frac{1}{8}}\right)}^2}{2}=\frac{2^{\frac{9}{4}}}{2}=2^{\frac{5}{4}} \Rightarrow D.$




\textbf{Solution 2 (Algebraic approach via Sum-to-Product Identity)}

Observing the coefficients of the given equation $z^4 + 4iz^3 - 6z^2 - 4iz - i = 0$, we notice its similarity to the binomial expansion of $(z+i)^4$. Recall that:
 \[(z+i)^4 = z^4 + 4iz^3 + 6(i)^2z^2 + 4(i)^3z + i^4 = z^4 + 4iz^3 - 6z^2 - 4iz + 1\]Thus, the original equation can be rewritten as:
 \[(z+i)^4 - 1 - i = 0 \implies (z+i)^4 = 1+i\]Let $w = z+i$. In polar form, $1+i = \sqrt{2} e^{i\frac{\pi}{4}}$. The four roots for $w$ are given by:
 \[w_k = \left(2^{\frac{1}{2}}\right)^{\frac{1}{4}} e^{i\left(\frac{\pi/4 + 2k\pi}{4}\right)} = 2^{\frac{1}{8}} e^{i\left(\frac{\pi}{16} + \frac{k\pi}{2}\right)}, \quad k \in \{0, 1, 2, 3\}\]Since $z_k = w_k - i$, the side length $s$ of the square formed by these vertices is the distance between two adjacent roots, such as $z_0$ and $z_1$. Note that the constant translation $-i$ cancels out when taking the difference:
 \[s = |z_1 - z_0| = |w_1 - w_0| = 2^{\frac{1}{8}} \left| e^{i\frac{9\pi}{16}} - e^{i\frac{\pi}{16}} \right|\]Expanding using Euler's formula $e^{i\theta} = \cos\theta + i\sin\theta$:
 \[e^{i\frac{9\pi}{16}} - e^{i\frac{\pi}{16}} = \left( \cos\frac{9\pi}{16} - \cos\frac{\pi}{16} \right) + i \left( \sin\frac{9\pi}{16} - \sin\frac{\pi}{16} \right)\]Applying the sum-to-product identities:

 
            \begin{itemize}
        \item $\cos A - \cos B = -2\sin\frac{A+B}{2}\sin\frac{A-B}{2}$
        \item $\sin A - \sin B = 2\cos\frac{A+B}{2}\sin\frac{A-B}{2}$
            \end{itemize}
            Setting $A = \frac{9\pi}{16}$ and $B = \frac{\pi}{16}$, we have $\frac{A+B}{2} = \frac{5\pi}{16}$ and $\frac{A-B}{2} = \frac{\pi}{4}$:
 \[e^{i\frac{9\pi}{16}} - e^{i\frac{\pi}{16}} = -2\sin\frac{5\pi}{16}\sin\frac{\pi}{4} + i \left( 2\cos\frac{5\pi}{16}\sin\frac{\pi}{4} \right) = 2\sin\frac{\pi}{4} \left( -\sin\frac{5\pi}{16} + i\cos\frac{5\pi}{16} \right)\]The magnitude of the term in the parentheses is $\sqrt{(-\sin\theta)^2 + (\cos\theta)^2} = 1$. Given $\sin\frac{\pi}{4} = 2^{-\frac{1}{2}}$, the side length is:
 \[s = 2^{\frac{1}{8}} \cdot \left| 2 \cdot 2^{-\frac{1}{2}} \cdot 1 \right| = 2^{\frac{1}{8}} \cdot 2^{\frac{1}{2}} = 2^{\frac{5}{8}}\]The area of the square is $s^2$:
 \[\text{Area} = (2^{\frac{5}{8}})^2 = 2^{\frac{10}{8}} = 2^{\frac{5}{4}}\]The answer is $\text{(D)}$.

 


	\item (\textit{AMC 12 12A 2008}) Triangle $ABC$ has $\angle C = 60^{\circ}$ and $BC = 4$.  Point $D$ is the midpoint of $BC$.  What is the largest possible value of $\tan{\angle BAD}$?


 $\mathrm{(A)}\ \frac{\sqrt{3}}{6}\qquad\mathrm{(B)}\ \frac{\sqrt{3}}{3}\qquad\mathrm{(C)}\ \frac{\sqrt{3}}{2\sqrt{2}}\qquad\mathrm{(D)}\ \frac{\sqrt{3}}{4\sqrt{2}-3}\qquad\mathrm{(E)}\ 1$

 \textbf{Answer: }D



\textbf{Solution 1}


                    \begin{figure}[H]
                        \centering
                        \includegraphics[width=0.5\textwidth]{Solution/AMC_12_12A_2008_24_Solution_Figure_1.png}
                    \end{figure}
                    

 Let $x = CA$. Then $\tan\theta = \tan(\angle BAF - \angle DAE)$, and since $\tan\angle BAF = \frac{2\sqrt{3}}{x-2}$ and $\tan\angle DAE = \frac{\sqrt{3}}{x-1}$, we have

 
 \[\tan\theta = \frac{\frac{2\sqrt{3}}{x-2} - \frac{\sqrt{3}}{x-1}}{1 + \frac{2\sqrt{3}}{x-2}\cdot\frac{\sqrt{3}}{x-1}}= \frac{x\sqrt{3}}{x^2-3x+8}\]

 With calculus, taking the \href{https://artofproblemsolving.com/wiki/index.php?title=Derivative}{derivative} and setting equal to zero will give the maximum value of $\tan \theta$. Otherwise, we can apply \href{https://artofproblemsolving.com/wiki/index.php?title=AM-GM}{AM-GM}:

 
 \begin{align*} \frac{x^2 - 3x + 8}{x} = \left(x + \frac{8}{x}\right) -3 &\geq 2\sqrt{x \cdot \frac 8x} - 3 = 4\sqrt{2} - 3\\ \frac{x}{x^2 - 3x + 8} &\leq \frac{1}{4\sqrt{2}-3}\\ \frac{x\sqrt{3}}{x^2 - 3x + 8} = \tan \theta &\leq \frac{\sqrt{3}}{4\sqrt{2}-3}\end{align*}


 Thus, the maximum is at  $\frac{\sqrt{3}}{4\sqrt{2}-3} \Rightarrow \mathbf{(D)}$.

 



\textbf{Solution 2}

We notice that $\tan(x)$ is strictly increasing on the interval $[0, \frac{\pi}{2})$ (if $\angle BAD\ge 90^\circ$, then it is impossible for $\angle C=60^\circ$), so we want to maximize $\angle BAD$. 

 Consider the circumcircle of $BAD$ and let it meet $AC$ again at $F$. Any point $P$ between $A$ and $F$ on line $AC$ is inside this circle, so it follows that $\angle BPD>\angle BAD$. Therefore to maximize $\angle BAD$, the circumcircle of $BAD$ must be tangent to $AC$ at $A$. By PoP we find that $CA^2=CD\cdot CB \Rightarrow AC = 2\sqrt{2}$.

 Now our computations are straightforward:
 \[\tan\angle BAD = \frac{\sin \angle BAD}{\cos \angle BAD} = \frac{\frac{2\sin\angle ABD}{AD}}{\frac{AB^2+AD^2-BD^2}{2AB\cdot AD}}\]
 \[=\frac{4\sin \angle ABD\cdot AB}{AB^2+AD^2-4} = \frac{4 AC \sin \angle ACB}{AB^2 + AD^2 - 4}\]
 \[=\frac{4\sqrt{6}}{(4^2+(2\sqrt{2})^2-4\cdot 2\sqrt{2}) + (2^2+(2\sqrt{2})^2 - 2\cdot 2\sqrt{2}) - 4} = \frac{4\sqrt{6}}{32-12\sqrt{2}}\]
 \[=\boxed{\frac{\sqrt{3}}{4\sqrt{2}-3}}\]



	\item (\textit{AMC 12 12A 2008}) A sequence $(a_1,b_1)$, $(a_2,b_2)$, $(a_3,b_3)$, $\ldots$ of points in the coordinate plane satisfies


 $(a_{n + 1}, b_{n + 1}) = (\sqrt {3}a_n - b_n, \sqrt {3}b_n + a_n)$  for $n = 1,2,3,\ldots$.


 Suppose that $(a_{100},b_{100}) = (2,4)$.  What is $a_1 + b_1$?


 $\mathrm{(A)}\ -\frac{1}{2^{97}}\qquad\mathrm{(B)}\ -\frac{1}{2^{99}}\qquad\mathrm{(C)}\ 0\qquad\mathrm{(D)}\ \frac{1}{2^{98}}\qquad\mathrm{(E)}\ \frac{1}{2^{96}}$

 \textbf{Answer: }D



\textbf{Solution 1}

This sequence can also be expressed using matrix multiplication as follows: 

 $\left[ \begin{array}{c} a_{n+1} \\ b_{n+1} \end{array} \right] = \left[ \begin{array}{cc} \sqrt{3} & -1 \\ 1 & \sqrt{3} \end{array} \right] \left[ \begin{array}{c} a_{n} \\ b_{n} \end{array} \right] = 2 \left[ \begin{array}{cc} \cos 30^\circ & -\sin 30^\circ \\ \sin 30^\circ & \ \cos 30^\circ \end{array} \right] \left[ \begin{array}{c} a_{n} \\ b_{n} \end{array} \right]$. 

 Thus, $(a_{n+1} , b_{n+1})$ is formed by rotating $(a_n , b_n)$ counter-clockwise about the origin by $30^\circ$ and dilating the point's position with respect to the origin by a factor of $2$. 

 So, starting with $(a_{100},b_{100})$ and performing the above operations $99$ times in reverse yields $(a_1,b_1)$. 

 Rotating $(2,4)$ clockwise by $99 \cdot 30^\circ \equiv 90^\circ$ yields $(4,-2)$. A dilation by a factor of $\frac{1}{2^{99}}$ yields the point $(a_1,b_1) = \left(\frac{4}{2^{99}}, -\frac{2}{2^{99}} \right) = \left(\frac{1}{2^{97}}, -\frac{1}{2^{98}} \right)$. 

 Therefore, $a_1 + b_1 = \frac{1}{2^{97}} - \frac{1}{2^{98}} = \frac{1}{2^{98}} \Rightarrow D$.

 



\textbf{Solution 2 (algebra)}

Let $(x,y)=(a_1,b_1)$. Then, we can begin to list out terms as follows:

 


 $(a_2,b_2)=(x\sqrt{3}-y,y\sqrt{3}+x)$
$(a_3,b_3)=(2x-2y\sqrt{3},2y+2x\sqrt{3})$
$(a_4,b_4)=(-8y,8x)$

 


 We notice that the sequence follows the rule $(a_{n+3},b_{n+3})=(-2^3b_n,2^3a_n)$

 We can now start listing out every third point, getting:

 


 $(a_1,b_1)=(x,y)$
$(a_4,b_4)=(-2^3y,2^3x)$
$(a_7,b_7)=(-2^6x,-2^6y)$
$(a_{10},b_{10})=(2^9y,-2^9x)$
$(a_{13},b_{13})=(2^{12}x,2^{12}y)$

 


 We can make two observations from this: 

 (1) In $a_n$, the coefficient of $x$ and $y$ is $2^{n-1}$

 (2) The positioning of $x$ and $y$, and their signs, cycle with every $12$ terms.

 


 We know then that from (1), the coefficients of $x$ and $y$ in $(a_{100},b_{100})$ are both $2^{99}$

 We can apply (2), finding $100 \text{(mod }12)=4$, so the positions and signs of $x$ and $y$ are the same in $(a_{100},b_{100})$ as they are in $(a_{4},b_{4})$.

 From this, we can get $(a_{100},b_{100})=(-2^{99}y,2^{99}x)$. We know that $(a_{100},b_{100})=(2,4)$, so we get the following:

 


 $-2^{99}y=2 \Rightarrow y=-\frac{1}{2^{98}}$
$2^{99}x=4 \Rightarrow x=\frac{1}{2^{97}}$

 


 The answer is $x+y=\frac{1}{2^{97}}-\frac{1}{2^{98}}=\boxed{\textbf{(D) }\frac{1}{2^{98}}}$..

 



\textbf{Solution 3}

The ordered pairs and $\sqrt{3}$'s makes us think to use complex numbers. We have $(a_{n+1},b_{n+1}) = 2\left(\frac{\sqrt{3}}{2}a_n - \frac{1}{2}b_n, \frac{\sqrt{3}}{2}b_n + \frac{1}{2}a_n\right)$, so $a_{n+1} + b_{n+1}i = 2\left(\frac{\sqrt{3}}{2} + \frac{1}{2}i\right)(a_n + b_ni) = \frac{1}{2}e^{i\pi/6}(a_n + b_ni)$. Letting $z_n = a_n + b_ni$ (so $z_{n+1} = a_{n+1} + b_{n+1}i$), we have $z_{n+1} = 2e^{i\pi/6}z_n$. Letting $n\rightarrow n-1$, we have $z_n = 2e^{i\pi/6}z_{n-1}$, so $z_{n-1} = \frac{1}{2}e^{-i\pi/6}z_n$. This is the reverse transformation. We have
 \[z_{99} = \frac{1}{2}e^{-i\pi/6}z_{100}\]
 \[z_{98} = \frac{1}{2^2}e^{2(-i\pi/6)}z_{100}\]
 \[\vdots\]
 \[z_{1} = \frac{1}{2^{99}}e^{99(-i\pi/6)}z_{100}\]
 \[= \frac{1}{2^{99}}e^{-i\pi/2}z_{100} = -\frac{1}{2^{99}}i(2 + 4i) = \frac{1}{2^{97}} - \frac{1}{2^{98}}i.\]

 Hence, $a_1 + b_1 = \frac{1}{2^{97}} - \frac{1}{2^{98}} = \boxed{\mathbf{(D)}\frac{1}{2^{98}}}$ ~ brainfertilzer.

 


 



\textbf{Solution 4 (Kinda braindead)}

Start by turning the two equations into $a_n = \frac{a_{n+1}\sqrt{3}+b_{n+1}}{4}$ and $b_n = \frac{b_{n+1}\sqrt{3}-a_{n+1}}{4}$. Note that these are just obtained by solving the equations. 

 
This makes finding values of $a_n$ and $b_n$ much easier, and soon we notice that $a_{97} = \frac{1}{2}$ and $b_{97} = -\frac{1}{4}$. After that, we get that $a_{94} = -\frac{1}{32}$ and $b_{94} = -\frac{1}{16}$. Observe that $|a_n| = |\frac{b_{n+3}}{8}|$ and $|b_n| = |\frac{a_{n+3}}{8}|$. This is basically just ignoring signs.

 
Now, we proceed to find $|a_1| = \frac{1}{2^{97}}$ while $|b_1| = \frac{1}{2^{98}}$. Despite there being 4 possible sign combinations for $(a_1, b_1)$, the only achievable answer choice is $\boxed{\textbf{(D)}\frac{1}{2^{98}}}$

 
-skibbysiggy

 


\end{enumerate}